% !TEX root = ../article.tex
% =============================================================================
% 3.2 Resumen ejecutivo
% =============================================================================
% La guía pide: descripción del problema, la solución desarrollada, los datos
% utilizados y los principales resultados. Debe permitir comprender el aporte
% del proyecto y NO EXCEDER UNA PÁGINA. Si se pasa, los scripts de
% compilación lo avisan (ver \articleCheckOnePage en config/commands.tex).
%
% Versión del primer avance: resume lo redactado en las secciones 1 a 3. Los
% resultados se completan cuando existan y el resto se ajusta tras cada
% asesoría.
% =============================================================================

\section*{Resumen ejecutivo}
\addcontentsline{toc}{section}{Resumen ejecutivo}
\label{summary:start}

\paragraph{Problema.}
Las enfermedades pulmonares están entre las principales causas de muerte en el
mundo: en 2023, la EPOC fue la tercera, y en 2024 el cáncer de pulmón fue el
más diagnosticado, con casi 2,6 millones de casos nuevos
\parencite{oms2026, sung2026}. La tomografía computarizada permite medirlas,
pero para cuantificar el volumen pulmonar o la extensión de una lesión primero
hay que delimitar el pulmón. Hacerlo a mano toma horas por estudio, los métodos
clásicos fallan cuando el pulmón tiene patología, y los modelos publicados son
difíciles de comparar y de reproducir.

\paragraph{Solución.}
Thoraxis segmenta automáticamente el pulmón izquierdo y el derecho en
tomografías de tórax con Deep Learning y calcula el volumen de cada uno. La
solución base es una U-Net 2D que procesa cada corte por separado; la propuesta
agrega contexto entre cortes contiguos (2.5D), aumentos de datos y datos de
entrenamiento más diversos. El modelo elegido se exporta a ONNX y se sirve con
una API en FastAPI, que alimenta un visor web con la máscara superpuesta y un
mapa de incertidumbre que muestra dónde el modelo es menos confiable. El
ciclo de vida sigue prácticas de MLOps: experimentos en MLflow sobre
Databricks, datos versionados con DVC, integración continua con GitHub Actions,
reentrenamiento en Kaggle, monitoreo de la deriva con Evidently y del
rendimiento de la API con OpenTelemetry y Grafana.

\paragraph{Datos.}
Se usan tres conjuntos públicos y anonimizados con máscaras de ambos pulmones.
Para entrenar, validar y probar se combinan TotalSegmentator, con 1939
tomografías de rutina clínica, y 20 tomografías de pacientes con COVID-19. La
validación externa se hace con LCTSC: 60 pacientes con cáncer de pulmón de tres
instituciones, que no se usan para entrenar. Toda partición se hace
por paciente, y ningún dato médico se guarda en los repositorios.

\paragraph{Resultados esperados.}
El proyecto se considerará exitoso si la propuesta alcanza un coeficiente de
Dice medio de al menos 0,95 por pulmón en la validación externa, supera a la
solución base con un intervalo de confianza que excluya el cero, funciona de
extremo a extremo en el prototipo y su evaluación puede reproducirse desde los
repositorios.
\pendiente{métricas obtenidas por la propuesta y la solución base}

\paragraph{Aporte y limitaciones.}
Thoraxis aporta un modelo de segmentación pulmonar portable, explicado con
mapas de incertidumbre en lugar de mapas de calor poco confiables, y una
evaluación reproducible, con validación externa e intervalos de confianza. Es
una
herramienta de investigación y no un sistema de diagnóstico: se evalúa solo con
datos públicos, sin validación clínica, y las diferencias en cómo cada conjunto
delimita el pulmón pueden afectar sus métricas.

\medskip
\noindent\textbf{Palabras clave:} \articleKeywords.

\label{summary:end}
\articleCheckOnePage
